\section{Positive words and tensor amplification}
\label{gen:sec:generation}

For an invertible square matrix $A$, write $|A|$ for its width.
A \emph{monomial matrix} is a permutation matrix times an invertible
diagonal matrix.  A \emph{call to $A$ on $w$ coordinates}, where $w\geq
|A|$, applies $A$ to an ordered tuple of $|A|$ distinct coordinates and
fixes the others.  A \emph{matrix word} is a finite composition of such
calls and full-width monomial matrices.  In a word using one specified
matrix $A$, its \emph{call count} counts only occurrences of $A$;
monomials have zero call count.  This convention concerns words alone.
In the scalar circuit model, a monomial on $w$ coordinates costs at most
$w$ scalar multiplications, with its permutation realized by designating
the appropriate available values.

\begin{definition}[Finite win]\label{def:finite-win}
A finite win consists of an invertible nonmonomial matrix
$A\in\GL_q(\C)$, an integer $b\geq2$, and a word on exactly $q^b$
coordinates implementing $A^{\otimes b}$ with $g$ calls to $A$, where
\begin{equation}\label{gen:eq:strict-win}
g<bq^{b-1}.
\end{equation}
\end{definition}

The comparison in \cref{gen:eq:strict-win} is with the usual tensor-axis
algorithm: for each of the $b$ axes, it makes $q^{b-1}$ calls to $A$.
Nonmonomiality implies $q\geq2$.  It also implies that $A^{\otimes b}$
is nonmonomial: an invertible nonmonomial matrix has a row with at least
two nonzero entries, and tensoring that row with any nonzero rows
preserves this property.  Thus a winning word necessarily has $g\geq1$.

\subsection{One fixed positive pattern generates every target}

The following statement supplies both the local substitutions in the
Fourier transfer and positive-call identity words used later in packing.
The positions of all calls are fixed before the target matrix is chosen.

\begin{samepage}
\begin{lemma}[Positive generation]\label{lem:generation}
Let $A\in\GL_q(\C)$ be nonmonomial.  There is an integer $h>0$ such that
every $H\in\GL_q(\C)$ admits a factorization
\begin{equation}\label{gen:eq:common-pattern}
H=M_0 A M_1 A\cdots A M_h,
\end{equation}
with exactly $h$ occurrences of $A$ and monomial matrices $M_i$.
In particular, identity admits such a factorization with a positive
number of calls.  More generally, for every fixed $w\geq q$, there is
a fixed positive pattern on $w$ coordinates realizing every element
of $\GL_w(\C)$ by embedded calls to $A$ and monomials.
\end{lemma}
\end{samepage}

\begin{proof}
We build a fixed word whose adjustable diagonal factors give a full-rank
differential, then show that its actual values contain a nonempty Zariski
open subset of $\GL_q(\C)$.  Every invertible target is a product of two
values from this subset, so doubling the word gives the common pattern.

Let $E_{ij}$ denote the matrix units.  Let
$\mathfrak d\subset\Mat_q(\C)$ be the diagonal subspace and let
$\mathcal S$ be the semigroup of words in $A$ and monomials, including
identity.  Inverses are not among the specified calls.  First we prove
\begin{equation}\label{gen:eq:tangent-span}
\Span_{s\in\mathcal S}s\mathfrak d s^{-1}=\Mat_q(\C).
\end{equation}
If $A\mathfrak d A^{-1}\subset\mathfrak d$, each rank-one idempotent
$A E_{ii}A^{-1}$ would be a diagonal rank-one idempotent, hence some
$E_{\pi(i)\pi(i)}$.  Orthogonality of these idempotents forces $\pi$
to be a permutation.  The identity
$A E_{ii}=E_{\pi(i)\pi(i)}A$ then says that column $i$ of $A$ is
supported in row $\pi(i)$; invertibility would make $A$ monomial.
Consequently some $X\in A\mathfrak d A^{-1}$ has $X_{ij}\ne0$ for
$i\ne j$.

Write $\mathcal V$ for the span on the left of
\cref{gen:eq:tangent-span}.  For every invertible diagonal
$D=\diag(d_1,\ldots,d_q)$, the matrix $DXD^{-1}$ belongs to $\mathcal V$.
If a linear functional $\ell$ vanishes on $\mathcal V$, then
\[
0=\ell(DXD^{-1})
 =\sum_a X_{aa}\ell(E_{aa})
  +\sum_{a\ne c}X_{ac}\ell(E_{ac})d_a/d_c
\qquad(d_1,\ldots,d_q\in\C^*).
\]
The Laurent monomials $1$ and $d_a/d_c$ for $a\ne c$ are distinct and
linearly independent.  Thus $X_{ij}\ell(E_{ij})=0$, and so
$\ell(E_{ij})=0$.  Finite-dimensional linear duality gives
$E_{ij}\in\mathcal V$.  Conjugation by any permutation preserves
$\mathcal V$, because left multiplication of a word by that permutation
is again a word.  Hence $\mathcal V$ contains every off-diagonal matrix
unit.  It also contains $\mathfrak d$, using $s=I_q$.  This proves
\cref{gen:eq:tangent-span}.

We next arrange the spanning directions along one word.  Start with
the prefix $h_0=I_q$ and its diagonal directions.  If the directions
collected so far span a proper subspace and the current prefix is $h$,
then
\[
\Span_{s\in\mathcal S}(hs)\mathfrak d(hs)^{-1}
 =h\Mat_q(\C)h^{-1}=\Mat_q(\C).
\]
Choose a word $s$ whose directions enlarge the collected span, and
append it to the right of $h$ in written product order.  Each choice
increases the dimension, so finitely many choices suffice.  Denote the
successive appended words by $s_1,\ldots,s_r$ and their prefixes by
$h_i=s_1\cdots s_i$.  The spaces $h_i\mathfrak d h_i^{-1}$, including
$i=0$, now span $\Mat_q(\C)$.

Insert variable diagonal factors in the product
\[
\Phi(D_0,\ldots,D_r)=D_0s_1D_1s_2\cdots s_rD_r.
\]
At $D_0=\cdots=D_r=I_q$, its value is $h_r$.  Varying $D_i$ in the
diagonal direction $E_{aa}$ gives the right-translated derivative
\[
(\delta\Phi)h_r^{-1}=h_iE_{aa}h_i^{-1}.
\]
Thus the derivative has rank $a=q^2$.  Choose $a$ diagonal entries
whose derivative columns are independent, and fix all other entries
to $1$.  Listing matrix entries produces a polynomial map
\[
f:\C^a\longrightarrow\C^a
\]
whose Jacobian determinant is nonzero at $(1,\ldots,1)$.  Parameters in
$(\C^*)^a$ give admissible words.  We show next that their image contains
an actual nonempty Zariski open subset of $\GL_q(\C)$.

Write $x_1,\ldots,x_a$ for the selected parameters and
$f_1,\ldots,f_a$ for the output entries.  The $f_i$ are algebraically
independent.  Indeed, a nonzero polynomial relation $P(f)=0$ of least
total degree would, upon differentiation and inversion of the Jacobian
over $\C(x_1,\ldots,x_a)$, imply
$(\partial P/\partial y_i)(f)=0$ for every $i$.  In characteristic zero
at least one of these is a nonzero polynomial of smaller degree,
a contradiction.  Both fields in the inclusion
\[
\C(f_1,\ldots,f_a)\subset\C(x_1,\ldots,x_a)
\]
therefore have transcendence degree $a$.  The latter is an algebraic
extension generated by the finitely many $x_i$, so it is finite.
Each of the finitely many
elements $x_i,x_i^{-1}$ satisfies a monic equation over the smaller
field.  Choose a nonzero polynomial $g(y_1,\ldots,y_a)$ divisible by
the denominators of all coefficients of these equations.  With
\begin{align*}
R&=\C[f_1,\ldots,f_a,1/g(f)],\\
B&=R[x_1,\ldots,x_a,x_1^{-1},\ldots,x_a^{-1}]
   \subset\C(x_1,\ldots,x_a),
\end{align*}
all the displayed generators of $B$ are integral over $R$.
Reducing their powers by their monic equations shows directly that
$B$ is a finite $R$-module.  Both are domains and $R\hookrightarrow B$
is injective.  This is the finite-type integral-extension criterion
\citep[Tag~02JJ, Lemma~10.36.5]{stacksProject}.

Fix any target $y\in\C^a$ with $g(y)\ne0$, and let
$\mathfrak m=(f_1-y_1,\ldots,f_a-y_a)\subset R$.  Algebraic independence
identifies $R/\mathfrak m$ with $\C$.  Localizing the inclusion above
gives $R_{\mathfrak m}\hookrightarrow B_{\mathfrak m}$ inside the field
$\C(x_1,\ldots,x_a)$, so $B_{\mathfrak m}$ is a nonzero finite module
over the local ring $R_{\mathfrak m}$.  Its residue algebra
\begin{equation}\label{gen:eq:nonzero-fiber}
B_{\mathfrak m}/\mathfrak mB_{\mathfrak m}
\end{equation}
is nonzero.  This is the lying-over step for the integral inclusion
\citep[Tag~00GQ, Lemma~10.36.17]{stacksProject}.
For completeness, if it were zero, a finite list of module
generators $b_1,\ldots,b_e$ would satisfy
$b_i=\sum_j c_{ij}b_j$ with all $c_{ij}$ in the maximal ideal of
$R_{\mathfrak m}$.  The determinant of $I-(c_{ij})$ is a unit in this
local ring.  Multiplication by its adjugate forces every $b_i$ to be
zero, contradicting $B_{\mathfrak m}\ne0$.

The algebra in \cref{gen:eq:nonzero-fiber} is finite dimensional over
$\C$.  A quotient by a maximal ideal is a finite field extension of
$\C$, and hence is $\C$ itself.  The resulting homomorphism
$B\to\C$ sends $x_i$ to some $\xi_i\ne0$, since $x_i^{-1}$ is also in
$B$.  The polynomial identities defining the $f_i$ descend to
$f_i(\xi)=y_i$.  These are actual admissible parameters of the fixed
word.  In particular its image contains
\[
\Omega=\{Y\in\Mat_q(\C):g(Y)\det Y\ne0\},
\]
a nonempty open set, because $g(Y)\det Y$ is a nonzero polynomial.

Finally fix $H\in\GL_q(\C)$.  There is a $V$ for which both
$V\in\Omega$ and $V^{-1}H\in\Omega$.  To see this without any limiting
argument, clear powers of $\det V$ from $g(V^{-1}H)$.  Its numerator is
a nonzero polynomial: the substitution $V\mapsto V^{-1}H$ is a
bijection of $\GL_q(\C)$, with inverse $W\mapsto HW^{-1}$, so this
rational function is not identically zero.  Avoid simultaneously the
zeros of that numerator, $g(V)$, and $\det V$.  A finite product of
nonzero polynomials over $\C$ is nonzero and has a point where it does
not vanish.  This supplies the required $V$.  Setting $W=V^{-1}H$
gives $H=VW$, with both factors realized by the same fixed word
pattern.  Two copies of that pattern therefore realize every $H$.

The initial pattern contains at least one occurrence of $A$: otherwise
its values have one fixed monomial support, a space of dimension $q$,
contradicting the derivative rank $q^2>q$.  Expanding the fixed words
$s_i$ and merging adjacent monomials in the doubled pattern gives
\cref{gen:eq:common-pattern} with a fixed $h>0$.  All inverses used in
the argument occurred in matrix identities or parameter selection;
every factor of the constructed words is a forward call to $A$ or a
monomial.  For $w>q$, applying the same result to the nonmonomial matrix
$A\oplus I_{w-q}$ proves the wider-width assertion, because each call
to this matrix is one embedded call to $A$.
\end{proof}

\subsection{Amplifying a finite saving}

The next lemma concerns scalar gates, including the cost of every
monomial scaling.  All its constants are allowed to depend on the
single fixed finite win.

\begin{lemma}[Tensor amplification]\label{lem:amplification}
If $(A,b,g)$ and a specified winning word form a finite win, then there
are constants $C<\infty$ and $0<\alpha<1$ such that $A^{\otimes k}$ has
an exact scalar linear circuit with at most
$Cq^k(k+1)^\alpha$ gates for every integer $k\geq0$.
\end{lemma}

\begin{proof}
Let $L_A(k)$ be the minimum scalar gate count for $A^{\otimes k}$,
and put $f(k)=L_A(k)/q^k$.  Here $A^{\otimes0}=I_1$ and $f(0)=0$.
These minima exist: entrywise multiplication and addition give finite
circuits, and a nonempty set of nonnegative integer gate counts has
a minimum.  Set
\[
Q=q^b,\qquad p=q/Q,\qquad j=\lfloor k/b\rfloor.
\]
Tensor the winning word with itself over $j$ groups, multiplying the
corresponding slots together.  The product of these slots is
$(A^{\otimes b})^{\otimes j}=A^{\otimes bj}$.

A monomial slot remains monomial after taking this tensor power, and
therefore costs at most $Q^j$ scalar gates.  Each call slot is
permutation-conjugate to $A\oplus I_{Q-q}$.  Its $j$-fold tensor power,
after reordering coordinates, has the following exact direct-sum
decomposition:
\begin{equation}\label{gen:eq:sector-decomposition}
(A\oplus I_{Q-q})^{\otimes j}
\ \cong\ 
\bigoplus_{r=0}^{j}
\left(A^{\otimes r}\right)^{\oplus
             \binom jr(Q-q)^{j-r}}.
\end{equation}
For a chosen subset of $r$ active factors, the inactive factors supply
$(Q-q)^{j-r}$ independent scalar coordinate choices, explaining the
multiplicity.  The total width is
$\sum_r\binom jr q^r(Q-q)^{j-r}=Q^j$.  Computing the blocks separately
with minimum circuits gives normalized cost at most
\begin{align}
\frac{1}{Q^j}\sum_{r=0}^j
  \binom jr(Q-q)^{j-r}L_A(r)
 &=\sum_{r=0}^j\binom jr p^r(1-p)^{j-r}f(r)\notag\\
 &=\mathbb E f(K),\qquad K\sim\operatorname{Bin}(j,p).
 \label{gen:eq:weighted-cost}
\end{align}
In particular, this distribution weights sectors by their coordinate
dimensions; the sectors have different widths.

Replicate the resulting circuit over the coordinates of the remaining
$k-bj<b$ tensor axes.  The normalized cost is unchanged by this
replication.  If a fixed circuit for $A$ has $\ell_A$ gates, applying
$A$ along each remaining axis costs $\ell_Aq^{k-1}$ gates per axis.
Since the number of monomial slots in the winning word is fixed,
there is a constant $C_0$ independent of $k$ such that
\begin{equation}\label{gen:eq:recurrence}
f(k)\leq C_0+g\,\mathbb E f(K),
\qquad K\sim\operatorname{Bin}(\lfloor k/b\rfloor,q/Q).
\end{equation}
For $j=0$ the grouped part is identity and the same inequality follows
from the bound on the fewer than $b$ remaining axes.  Every permutation
used in deriving \cref{gen:eq:sector-decomposition} is only a
reindexing; every diagonal scaling has already been charged.

Put $\rho=q/(Qb)$.  The strict saving gives $g\rho<1$, and hence we can
choose $0<\alpha<1$ sufficiently close to $1$ that
$g\rho^\alpha<1$.  Choose also a number $\gamma$ with
$g\rho^\alpha<\gamma<1$.  Since
\[
g\left(\frac{\rho k+1}{k+1}\right)^\alpha
 \longrightarrow g\rho^\alpha,
\]
there is $k_0\geq1$ such that this expression is at most $\gamma$ for all
$k\geq k_0$.  Choose $C$ at least $C_0/(1-\gamma)$ and large enough
that $f(k)\leq C(k+1)^\alpha$ for $k<k_0$.
For $k\geq k_0$, every value of $K$ satisfies
$K\leq\lfloor k/b\rfloor<k$.  Strong induction, concavity of
$x\mapsto x^\alpha$, and $\mathbb E K=p\lfloor k/b\rfloor\leq\rho k$
therefore give
\[
f(k)\leq C_0+gC\,\mathbb E(K+1)^\alpha
 \leq C_0+gC(\rho k+1)^\alpha
 \leq C_0+C\gamma(k+1)^\alpha
 \leq C(k+1)^\alpha.
\]
All circuits used in this induction compute exact matrices; the
expectation is merely the finite weighted sum in
\cref{gen:eq:weighted-cost}.  Multiplying by $q^k$ proves the scalar
gate bound
\[
L_A(k)=O\bigl(q^k(k+1)^\alpha\bigr),\qquad 0<\alpha<1.
\]
\end{proof}
